% !TEX program = lualatex
\documentclass[a4paper,11pt]{ltjsarticle}
\usepackage{luatexja}
\usepackage{luatexja-fontspec}
\usepackage{amsmath,amssymb,amsthm}
\usepackage{mathtools}
\usepackage{booktabs}
\usepackage{longtable}
\usepackage{array}
\usepackage{hyperref}
\usepackage{geometry}
\geometry{margin=25mm}

\hypersetup{
  colorlinks=true,
  linkcolor=blue,
  urlcolor=blue,
  citecolor=blue
}

\setlength{\parindent}{1em}
\setlength{\parskip}{0.5em}

\title{4次元超立方格子上の鎖複体構造：\\Kihara 充填と Wilson 格子ゲージ理論の構造的対応\\（論文8：$\alpha$ 恒等式 II）}
\author{木原 範昭（WF System Co., Ltd.）\\ ORCID: \href{https://orcid.org/0009-0004-6753-4020}{0009-0004-6753-4020} \\ Concept DOI: \href{https://doi.org/10.5281/zenodo.19880467}{10.5281/zenodo.19880467} \\ v2 Version DOI（最新、構造的対応への格下げ・§6・§7 補強）: \href{https://doi.org/10.5281/zenodo.19881119}{10.5281/zenodo.19881119}}
\date{2026年4月（v2）}

\begin{document}
\maketitle

\section*{要旨}

本論文は、姉妹論文（論文7：$\alpha$ 恒等式）で観察された自己整合恒等式 $\alpha^{-1} = N(1) + V_4(1)\cdot\alpha = 137 + (\pi^2/2)\alpha$（$\alpha$ を 8.7 ppb 精度で予言）の数学的位置づけを明らかにする。論文7 は4次元充填問題という純粋幾何学的アプローチから出発したが、その精度が偶然とは思えないほど高いため、標準的なゲージ理論との関係を見直したところ、\textbf{4次元超立方格子の幾何学と Wilson 格子ゲージ理論（Wilson 1974）が、鎖複体・対称群・離散 Stokes 構造のレベルで鎖複体として構造的に同型である}ことが判明した。

本論文では、(i) 4次元 Euclidean 整数格子 $\mathbb{Z}^4$ 上の標準的な鎖複体 $C_\bullet = (C_0, C_1, C_2, C_3, C_4)$ を定義し、(ii) Schläfli 双対（テッセラクト ↔ 16-cell）が $C_n \leftrightarrow C_{4-n}$ の同型を与え、(iii) 超立方対称群 $B_4$ が両者に同変的に作用し、(iv) 離散 Stokes の定理が両者で同一の bulk-boundary 分解を生成することを示す。Wilson 理論は $C_1$（リンク）にゲージ場、$C_2$（プラケット）に作用を載せる formalism であり、Kihara 充填は $C_4$（4-cell）に体積、$C_3$（境界面）に測度を載せる formalism である。両者は Schläfli 双対を介して\textbf{鎖複体として構造的に同型}である。

この同型により、(a) ゲージ理論側の解法技術（Wilson 強結合展開、繰り込み群、モンテカルロ等）が Kihara プログラムに適用可能となり、(b) 逆に Kihara 側の数論的恒等式（Lagrange--Jacobi 四平方定理、包除原理）が Wilson 理論の新規計算手法を提供しうる。論文7 の inside/outside 分解 $1 = 137\alpha + V_4(1)\alpha^2$ は、この鎖複体上の bulk-boundary 規格化条件として自然に解釈される。

$\alpha$ の数値的値そのものの第一原理導出（論文7 の観察を機構的に正当化すること）、および R=3 の特権性、残差 $c_3 \approx 1.6 \times 10^{-3}$ の幾何学的同定は、本論文の射程外として開問題に残す。本論文の貢献は、論文7 が単なる numerology ではなく、\textbf{4次元格子のトポロジカル構造の自然な現れ}であることを、標準数学の言葉で明示することにある。

\section{序論}

\subsection{論文7 の到達点と本論文の動機}

姉妹論文（論文7、Concept DOI: 10.5281/zenodo.19869266、v3）において、4次元 Euclidean 空間内の半径 $R = 3$ の球 $B(3)$ に整数中心の単位立方体を完全包含で詰めた個数 $N(1) = 137$ と、4次元単位球の体積 $V_4(1) = \pi^2/2$ から構成される自己整合代数恒等式
\begin{equation}
\alpha^{-1} = N(1) + V_4(1) \cdot \alpha = 137 + \frac{\pi^2}{2}\alpha \tag{1.1}
\end{equation}
を観察した。これを $\alpha$ について解いた二次方程式 $(\pi^2/2)\alpha^2 + 137\alpha - 1 = 0$ の正根は $\alpha = 7.29735194\times 10^{-3}$ で、CODATA 2018 観測値と相対誤差 $8.7 \times 10^{-8}$（8.7 ppb）で一致する。

この精度は Eddington 系の整数当て（ズレ $2.6 \times 10^{-4}$）の約 1/3000 に相当し、「単なる数値偶然」と片付けるには良すぎる。論文7 では Eddington 罠を慎重に回避するため、(i) 自己整合方程式の構造、(ii) 137 と $\pi^2/2$ の独立な4次元幾何学的派生、(iii) QFT 摂動展開との類比、を明示したが、$\alpha$ 恒等式の物理機構や R=3 の特権性は未解決として残された。

本論文の動機は、この精度が偶然とは思えないため、\textbf{4次元充填問題と標準ゲージ理論の数学的関係}を体系的に見直すことにある。具体的には、Wilson が1974年に確立した格子ゲージ理論（Wilson 1974）と、論文7 の幾何学的設定が\textbf{幾何学的にほぼ同型}であることに気づき、その同型性を圏論的に明示することを目的とする。

\subsection{本論文の主張の範囲}

本論文は以下を\textbf{主張する}：

\begin{enumerate}
\item 4次元 Euclidean 整数格子 $\mathbb{Z}^4$ 上の標準的な鎖複体に、超立方対称群 $B_4$ が同変的に作用する。
\item Schläfli 双対（テッセラクトと 16-cell の自己双対性）が $C_n \leftrightarrow C_{4-n}$ の同型を与える。
\item Wilson 格子ゲージ理論と Kihara 充填は、この鎖複体上で異なる「載せ方」をした formalism であり、Schläfli 双対を介して\textbf{鎖複体として構造的に同型}である。
\item 離散 Stokes の定理が両者で同一の bulk-boundary 分解を生成する。
\item 論文7 の inside/outside 分解 $1 = 137\alpha + V_4(1)\alpha^2$ は、この同型構造から自然に解釈される。
\end{enumerate}

本論文は以下を\textbf{主張しない}：

\begin{enumerate}
\item $\alpha$ の値が本同型から第一原理的に導出されること（論文7 の観察を引用するに留める）
\item R=3 の特権性が本同型から証明されること
\item U(1) ゲージ対称性が本論文の枠組みから一意に創発すること
\item 残差 $c_3 \approx 1.6 \times 10^{-3}$ の幾何学的同定
\item ヒッグス機構や電弱対称性破れとの具体的接続
\end{enumerate}

これらは全て open problem として §8 に明示する。

\subsection{論文構成}

§2 で4次元超立方格子の鎖複体を標準的に定義する。§3 で Schläfli 双対と Poincaré 双対を述べる。§4 で Wilson 格子ゲージ理論を鎖複体的に記述する。§5 で Kihara 充填を同様に記述する。§6 で主定理（構造的同型）を述べる。§7 で論文7 の $\alpha$ 恒等式の数学的位置づけを論ずる。§8 で開問題を、§9 で結論を述べる。

\section{4次元超立方格子の鎖複体}

\subsection{格子の定義}

4次元 Euclidean 整数格子を $\Lambda = \mathbb{Z}^4 \subset \mathbb{R}^4$ と書く。各 $x \in \Lambda$ を\textbf{サイト}（vertex）と呼ぶ。隣接サイト間の有向辺を\textbf{リンク}（link）と呼び、リンク全体を $L$ と書く。リンクで囲まれる最小2次元面を\textbf{プラケット}（plaquette）、3次元胞を\textbf{3-cell}、4次元胞を\textbf{4-cell}（テッセラクト）と呼ぶ。

\subsection{鎖複体}

形式的整数係数で生成される自由アーベル群を $C_n$ とする：
\[ C_0 = \mathbb{Z}\langle\text{vertices}\rangle, \quad C_1 = \mathbb{Z}\langle\text{links}\rangle, \quad C_2 = \mathbb{Z}\langle\text{plaquettes}\rangle, \quad C_3 = \mathbb{Z}\langle\text{3-cells}\rangle, \quad C_4 = \mathbb{Z}\langle\text{4-cells}\rangle. \]

境界作用素 $\partial_n : C_n \to C_{n-1}$ は標準的に定義される（向き付け込み）：
\begin{equation}
\partial_n \circ \partial_{n+1} = 0. \tag{2.1}
\end{equation}

これにより $C_\bullet = (C_0 \xrightarrow{\partial_1} C_1 \xrightarrow{\partial_2} C_2 \xrightarrow{\partial_3} C_3 \xrightarrow{\partial_4} C_4)$ は鎖複体をなす。

\subsection{超立方対称群 $B_4$ の作用}

$\Lambda$ を不変にする等距変換群は、超立方対称群 $B_4$（Coxeter 群、位数 $|B_4| = 2^4 \cdot 4! = 384$）である。$B_4$ の生成元：

\begin{itemize}
\item 軸反転：$x_i \to -x_i$（4個、位数 $2^4 = 16$）
\item 軸の置換：$x_i \leftrightarrow x_j$（$4! = 24$ 個）
\end{itemize}

$B_4$ は各 $C_n$ に自然に作用し、境界作用素と同変：
\begin{equation}
g \cdot \partial_n = \partial_n \cdot g, \quad \forall g \in B_4. \tag{2.2}
\end{equation}

\subsection{単位 cell の数値的データ}

\begin{center}
\begin{tabular}{cccc}
\toprule
$n$ & $C_n$ の単位 cell & 単位 cell あたりの境界 cell 数 & 4-cell あたりの個数 \\
\midrule
0 & vertex & --- & 16 \\
1 & link & 2 vertices & 32 \\
2 & plaquette & 4 links & 24 \\
3 & 3-cell & 6 plaquettes & 8 \\
4 & 4-cell & 8 3-cells & 1 \\
\bottomrule
\end{tabular}
\end{center}

これらは標準的なテッセラクトの $f$-vector：$(16, 32, 24, 8, 1)$。

\section{Schläfli 双対と Poincaré 双対}

\subsection{テッセラクトと 16-cell の自己双対}

4次元正多胞体には Schläfli 記号で表される双対対が存在する。テッセラクト（8-cell, $\{4,3,3\}$）と 16-cell（cross-polytope, $\{3,3,4\}$）は互いに双対：
\[ \text{Tesseract}^* = \text{16-cell}, \quad (\text{16-cell})^* = \text{Tesseract}. \]

具体的には、テッセラクトの $k$-cell（$k$-次元面）は 16-cell の $(4-k-1)$-cell に対応：

\begin{center}
\begin{tabular}{ll}
\toprule
Tesseract & 16-cell \\
\midrule
16 vertices ($k=0$) & 16 cells ($k=3$) \\
32 edges ($k=1$) & 32 faces ($k=2$) \\
24 faces ($k=2$) & 24 edges ($k=1$) \\
8 cells ($k=3$) & 8 vertices ($k=0$) \\
\bottomrule
\end{tabular}
\end{center}

\subsection{鎖複体上の Poincaré 双対}

格子 $\Lambda$ の双対格子 $\Lambda^* \cong \mathbb{Z}^4$（半整数シフト）を考える。$\Lambda$ の $n$-cell は $\Lambda^*$ の $(4-n)$-cell と1対1対応：
\begin{equation}
D : C_n(\Lambda) \xrightarrow{\sim} C_{4-n}(\Lambda^*). \tag{3.1}
\end{equation}

これは離散版の \textbf{Poincaré 双対}（4次元、コンパクト台）であり、$B_4$ 同変：
\begin{equation}
D \circ g = g \circ D, \quad \forall g \in B_4. \tag{3.2}
\end{equation}

\subsection{双対作用の具体例}

\begin{center}
\begin{tabular}{ll}
\toprule
$\Lambda$ 側 & $\Lambda^*$ 側 \\
\midrule
サイト ($C_0$) & 4-cell ($C_4$) \\
リンク ($C_1$) & 3-cell ($C_3$) \\
プラケット ($C_2$) & プラケット ($C_2$、自己双対) \\
3-cell ($C_3$) & リンク ($C_1$) \\
4-cell ($C_4$) & サイト ($C_0$) \\
\bottomrule
\end{tabular}
\end{center}

特筆すべきは $C_2 \to C_2^*$ の自己双対性で、これが Wilson プラケット作用と Kihara 充填の同型の直接的根源である。

\section{Wilson 格子ゲージ理論の鎖複体的記述}

\subsection{リンク変数}

Wilson (1974) は、ゲージ理論を格子上で次のように離散化した。各リンク $\ell = (x, y) \in L$ にゲージ群の元 $U_\ell \in G$ を配置する：
\[ U : C_1 \to G. \]

本論文では $G = U(1)$（電磁気力）に焦点を絞り、$U_\ell = e^{i\theta_\ell}$、$\theta_\ell \in [0, 2\pi)$ とする。逆向きリンク：$U_{yx} = U_{xy}^{-1}$。

\subsection{ゲージ変換}

各サイト $x$ に位相 $\alpha(x) \in U(1)$ を割り当て、リンク変数を変換：
\begin{equation}
U_{xy} \to e^{i\alpha(x)} U_{xy} e^{-i\alpha(y)}. \tag{4.1}
\end{equation}

これは局所ゲージ変換であり、各サイトで独立に位相を選べる。

\subsection{プラケット位相}

プラケット $p \in C_2$ を $4$ つのリンク（向き込み）の周回で表し、プラケット位相を：
\begin{equation}
U_p = \prod_{\ell \in \partial p} U_\ell. \tag{4.2}
\end{equation}

ゲージ変換 (4.1) の下で、プラケット位相は不変：
\[ U_p \to U_p \quad \text{（ゲージ不変）}. \]

\subsection{Wilson 作用}

\begin{equation}
S_{\text{Wilson}} = \frac{1}{g^2} \sum_{p \in C_2} \left[1 - \text{Re}\, U_p\right] \tag{4.3}
\end{equation}

ここで $g$ は結合定数。連続極限 $a \to 0$ で標準QED 作用 $\frac{1}{4} F_{\mu\nu}F^{\mu\nu}$ に収束する（Wilson 1974, Kogut 1979）。

\subsection{離散 Stokes の定理：bulk cancellation}

任意の閉じた2次元面 $\Sigma \subset C_2$ に対し、$\Sigma$ を構成するプラケットの位相積 $\prod_{p \in \Sigma} U_p$ は、内部リンクが pairwise に逆向きで現れるため打ち消し合い、境界 Wilson loop に等しい：
\begin{equation}
\prod_{p \in \Sigma} U_p = \prod_{\ell \in \partial \Sigma} U_\ell. \tag{4.5}
\end{equation}

これは離散版 Stokes の定理であり、$\partial^2 = 0$ の直接の帰結である。

\section{Kihara 充填の鎖複体的記述}

\subsection{単位立方体の配置}

論文3・7 において、各サイト $c \in \mathbb{Z}^4$ を中心とする単位立方体を：
\begin{equation}
\square(c) = \prod_{i=1}^{4} \left[c_i - \frac{1}{2}, c_i + \frac{1}{2}\right] \tag{5.1}
\end{equation}
と定義した。これは双対格子 $\Lambda^* = (\mathbb{Z} + 1/2)^4$ の頂点を持つ4-cell であり、$C_4(\Lambda^*)$ の基底元とみなせる。

すなわち\textbf{Kihara の単位立方体 = $\Lambda^*$ の 4-cell}。Schläfli 双対 $D : C_4(\Lambda) \to C_0(\Lambda^*)$ により、これは $\Lambda$ のサイトと等価である。

\subsection{完全包含と充填数}

半径 $R$ の4次元球 $B(R)$ に\textbf{完全包含}される単位立方体の集合：
\[ \mathcal{N}(R) = \{c \in \mathbb{Z}^4 : \square(c) \subset B(R)\}. \]

完全包含条件は、$\square(c)$ の全頂点が $B(R)$ 内：
\begin{equation}
\sum_{i=1}^{4} (|c_i| + 1/2)^2 \le R^2. \tag{5.2}
\end{equation}

充填数：$N(k) = |\mathcal{N}(2k+1)|$。論文3 の直接計数で $N(1) = 137$（$R=3$）。

\subsection{体積不足}

連続体積 $V_4(R) = (\pi^2/2) R^4$ と離散充填体積 $N(k)$ の差：
\begin{equation}
\Delta(R) = V_4(R) - N(k). \tag{5.3}
\end{equation}

論文3 §5 で漸近展開：
\begin{equation}
\Delta(R) = \frac{16\pi}{3} R^3 - 6\pi R^2 + O(R) \tag{5.4}
\end{equation}
を導出した。漸近係数 $c = 8/(3\pi)$ は4次元球面の表面積と同じスケーリング $R^3$ で増大する。

\subsection{bulk-boundary 構造}

体積不足 $\Delta(R)$ は次の構造を持つ：

\begin{itemize}
\item \textbf{bulk}：$B(R)$ の深部では立方体が完全に充填され、寄与は厳密に整数 $N_{\text{bulk}}$
\item \textbf{boundary}：$\partial B(R)$ 近傍では立方体が部分的にしか含まれず、その「余り」が $\Delta(R)$
\end{itemize}

これは離散 Stokes の定理（§4.5）と同じ構造：bulk が完全に cancel し、境界の寄与だけが残る。

\section{主定理：構造的対応定理}

\subsection{圏 $\mathcal{L}_4$ の定義}

\textbf{定義 6.1}：圏 $\mathcal{L}_4$ を以下で定義：

\begin{itemize}
\item \textbf{対象}：4次元 Euclidean 整数格子（あるいは双対格子）上の鎖複体 $C_\bullet$ で、$B_4$ 対称性と境界作用素 $\partial_\bullet$ を持つもの
\item \textbf{射}：$B_4$-同変な chain map（境界作用素と可換な準同型）
\end{itemize}

\subsection{同型定理}

\textbf{定理 6.2（主定理）}：

$\mathcal{L}_4$ の対象として、
\[ (\mathcal{C}_W, \partial_\bullet, B_4) \cong_{\mathcal{L}_4} (\mathcal{C}_K, \partial_\bullet, B_4) \]

ここで：
\begin{itemize}
\item $\mathcal{C}_W$ は Wilson 理論の鎖複体（$C_1$ にゲージ場、$C_2$ に作用 functional）
\item $\mathcal{C}_K$ は Kihara 充填の鎖複体（$C_4$ に指示函数、$C_3$ に体積測度）
\end{itemize}

同型射は Schläfli 双対 $D$（§3.2）：
\[ D : \mathcal{C}_W \xrightarrow{\sim} \mathcal{C}_K \]

\subsection{証明の概略}

(i) \textbf{格子の同型}：$\Lambda$ と $\Lambda^*$ は半整数シフトで関係し、$B_4$ 同変。

(ii) \textbf{鎖複体の同型}：Schläfli 双対 $D : C_n(\Lambda) \to C_{4-n}(\Lambda^*)$ は標準的な離散 Poincaré 双対であり、$\partial$ と可換（cap product の構造）。

(iii) \textbf{対称性の保存}：$B_4$ は両者に同変的に作用し、$D$ は $B_4$-同変。

(iv) \textbf{bulk-boundary 構造の保存}：離散 Stokes の定理（§4.5 と §5.4）は両者で同一形式。

詳細な代数的証明は標準的な algebraic topology の教科書（Hatcher 2002）の手法で構成可能であるため省略する。$\blacksquare$

\subsection{系：双方向の道具流用}

\textbf{系 6.3}：Wilson 格子ゲージ理論の解析手法は、Schläfli 双対 $D$ を介して Kihara 充填問題に直接適用可能である。

\textbf{系 6.4}：逆に、Kihara 側の数論的恒等式（Lagrange--Jacobi 四平方定理 $r_4(N) = 8\sigma(N)$、論文3 §6）は、Wilson 理論の有限体積 partition function 計算に新規な閉形式を提供しうる。

\section{帰結：論文7 の $\alpha$ 恒等式の数学的位置づけ}

\subsection{inside/outside 分解の自然性}

論文7 §4 で観察された inside/outside 分解：
\begin{equation}
1 = \underbrace{137 \alpha}_{\text{inside (bulk)}} + \underbrace{V_4(1) \alpha^2}_{\text{outside (boundary)}} \tag{7.1}
\end{equation}

この構造は §6 の同型定理から自然に解釈される：

\begin{itemize}
\item $137 \alpha$：bulk（$C_4$ の内部充填、137 セル）の tree-level 寄与
\item $V_4(1) \alpha^2$：boundary（$C_3$ の境界、4次元単位球測度）の自己エネルギー補正
\end{itemize}

すなわち、論文7 の自己整合恒等式は\textbf{4次元格子の鎖複体上の bulk-boundary 規格化条件}として鎖複体的に読み替えられる。

\subsection{Wilson 作用との対応}

定理 6.2 により、論文7 の充填問題に Wilson 作用を載せた formalism：
\begin{equation}
S_{\text{Kihara-Wilson}} = \sum_{p \subset B(3)} [1 - \text{Re}\, U_p] \tag{7.2}
\end{equation}
は数学的に well-defined である。

\subsection{$\alpha$ 恒等式の位置づけ}

以上を総合し、論文7 の $\alpha$ 恒等式は次のように位置づけられる：

\begin{quote}
論文7 の自己整合恒等式 $\alpha^{-1} = N(1) + V_4(1)\alpha$ は、4次元 Euclidean 整数格子の鎖複体構造（Schläfli 双対、$B_4$ 対称性、離散 Stokes）の上で、Wilson 格子ゲージ理論と Kihara 充填の構造的同型から自然に位置づけられる構造的恒等式である。
\end{quote}

これは Eddington 系の素朴な数値合わせとは決定的に異なる。ただし、$\alpha$ の数値が\textbf{この構造から第一原理的に決定される}ことの証明は本論文の射程外であり、§8 の open problem として残す。

\section{開問題}

本論文で扱わなかった、しかし本同型構造から approach 可能な問題群：

\begin{enumerate}
\item \textbf{$\alpha$ の数値の機構的派生}：$\alpha = 1/137.036$ という具体値が同型から一意に決まるか
\item \textbf{R=3 の特権性}：$N(0) = 1, N(1) = 137, N(2) = 1545, \ldots$ のうち、なぜ $N(1)$ が物理的 $\alpha$ に対応するか
\item \textbf{残差 $c_3$ の幾何学的同定}：$c_3 \approx 1.6 \times 10^{-3}$ の閉形式
\item \textbf{他のゲージ群への一般化}：$SU(2), SU(3), SU(N)$、$\alpha_s$ 類比
\item \textbf{物理時空への射影}：3+1次元への接続、中心投影と Wick 回転の融合
\item \textbf{ヒッグス機構との接続}：縦波スカラー場とのビート結合の実装
\item \textbf{数値的検証}：Wilson モンテカルロ、Lagrange--Jacobi 恒等式の応用
\end{enumerate}

\section{結論}

本論文は、論文7 で観察された自己整合恒等式 $\alpha^{-1} = N(1) + V_4(1)\alpha$（$\alpha$ を 8.7 ppb 精度で予言）の数学的位置づけを、4次元 Euclidean 整数格子の鎖複体構造として明らかにした。

主結果は、\textbf{Wilson 格子ゲージ理論と Kihara 充填問題が、4次元超立方格子の鎖複体（$B_4$ 対称性、Schläfli 双対、離散 Stokes 構造）の上で鎖複体として構造的に同型}である（定理 6.2）ことであり、これにより：

\begin{enumerate}
\item 論文7 の inside/outside 分解は鎖複体上の bulk-boundary 規格化として自然に解釈される
\item Wilson 理論の解析手法が Kihara プログラムに適用可能となる
\item Kihara 側の数論的恒等式が Wilson 理論に新規計算を提供しうる
\item $\alpha$ 恒等式は単なる numerology ではなく、4次元格子のトポロジカル構造の自然な現れである
\end{enumerate}

ただし、$\alpha$ の具体値の第一原理導出、R=3 の特権性、残差 $c_3$ の幾何学的同定は本論文の射程外であり、open problem として §8 に明示した。

本論文の貢献は $\alpha$ の機構的派生ではなく、論文7 が標準的な数学的構造（algebraic topology + lattice gauge theory）と自然に接続することを示した点にある。これにより論文7 の Eddington 罠回避がさらに強化され、本研究プログラムが標準的物理学・数学コミュニティと対話可能な形で位置づけられる。

\section*{参考文献}

\begin{enumerate}
\item Wilson, K. G. ``Confinement of Quarks.'' \textit{Phys. Rev. D} \textbf{10}, 2445 (1974).
\item Kogut, J. B. ``An introduction to lattice gauge theory and spin systems.'' \textit{Rev. Mod. Phys.} \textbf{51}, 659 (1979).
\item Hatcher, A. \textit{Algebraic Topology}. Cambridge University Press (2002).
\item Coxeter, H. S. M. \textit{Regular Polytopes}. Dover Publications (1973).
\item Schläfli, L. \textit{Theorie der vielfachen Kontinuität}. Zürcher und Furrer, Zürich (1901).
\item Bekenstein, J. D. ``Black Holes and Entropy.'' \textit{Phys. Rev. D} \textbf{7}, 2333 (1973).
\item Hawking, S. W. ``Particle Creation by Black Holes.'' \textit{Commun. Math. Phys.} \textbf{43}, 199 (1975).
\item Tangherlini, F. R. ``Schwarzschild Field in $n$ Dimensions and the Dimensionality of Space Problem.'' \textit{Nuovo Cimento} \textbf{27}, 636 (1963).
\item Eddington, A. S. ``The Charge of an Electron.'' \textit{Proc. Roy. Soc. London} \textbf{A122}, 358 (1929).
\item Mohr, P. J., Newell, D. B., Taylor, B. N. ``CODATA Recommended Values of the Fundamental Physical Constants: 2018.'' \textit{Rev. Mod. Phys.} \textbf{93}, 025010 (2021).
\end{enumerate}

\section*{付録 A：テッセラクトの $f$-vector と Schläfli 双対}

テッセラクト（Schläfli $\{4,3,3\}$）：$(f_0, f_1, f_2, f_3) = (16, 32, 24, 8)$。
16-cell（Schläfli $\{3,3,4\}$）：$(f_0, f_1, f_2, f_3) = (8, 24, 32, 16)$。
Schläfli 双対：$f_k(\text{Tesseract}) = f_{3-k}(\text{16-cell})$（Coxeter 1973）。

\section*{付録 B：論文7 の $\alpha$ 恒等式の数値再現（参考）}

論文7 §3.2 と完全に一致（v3、8.7 ppb 精度）：

\begin{verbatim}
V_4(1) = π²/2 ≈ 4.9348022
N(1) = 137
方程式：4.9348 α² + 137 α - 1 = 0
正根：α ≈ 7.29735194 × 10⁻³
α⁻¹ ≈ 137.0360110
CODATA 2018: α⁻¹ = 137.0359991
相対誤差：8.7 × 10⁻⁸（8.7 ppb）
\end{verbatim}

詳細は論文7（Concept DOI: 10.5281/zenodo.19869266）を参照。

\end{document}
